% 双环PID控制原理框图（PPT用，紧凑版）
% 编译: xelatex dual_loop_pid.tex
\documentclass[tikz,border=4pt]{standalone}
\usepackage[scheme=plain]{ctex}

\usetikzlibrary{shapes,arrows,positioning,calc,backgrounds,fit}

\begin{document}
\begin{tikzpicture}[
    auto,
    node distance=1.0cm and 1.0cm,
    >=stealth',
    block/.style={
        rectangle, draw=black, thick,
        text width=1.8cm, minimum height=0.8cm,
        align=center, font=\tiny\sffamily,
        rounded corners=2pt,
    },
    sum/.style={
        circle, draw=black, thick,
        minimum size=5mm, inner sep=0pt,
    },
    branch/.style={
        circle, fill=black,
        minimum size=1.2mm, inner sep=0pt,
    },
    dashbox/.style={
        draw=blue!50, dashed, thick,
        rounded corners=4pt, inner sep=5pt,
    },
    line/.style={draw=black, thick, ->},
    fline/.style={draw=black!55, thick, ->},
]

% ==== 输入 ====
\node [font=\tiny\sffamily] at (-1.6,0) {$r(t)$};
\node [sum, right=0.8cm] (s1) {};
\node at (s1) {\tiny $-$};
\draw [line] (-1.1,0) -- (s1);

% ==== 外环PID ====
\node [block, right=1.2cm of s1] (PID1) {外环PID};
\draw [line] (s1) -- (PID1);

% ==== 内环求和 ====
\node [sum, right=1.2cm of PID1] (s2) {};
\node at (s2) {\tiny $-$};
\draw [line] (PID1) -- (s2);

% ==== 内环PID ====
\node [block, right=1.2cm of s2] (PID2) {内环PID};
\draw [line] (s2) -- (PID2);

% ==== 被控对象 ====
\node [block, right=1.2cm of PID2, text width=2.0cm] (G2) {$G_2(s)$};
\node [block, right=1.0cm of G2, text width=2.0cm] (G1) {$G_1(s)$};
\draw [line] (PID2) -- (G2);
\draw [line] (G2) -- (G1);

% ==== 输出 ====
\draw [line] (G1.east) -- ++(0.8,0) node [right, font=\tiny\sffamily] {$y(t)$};
\coordinate (out) at ($(G1.east)+(0.4,0)$);
\node [branch] at (out) {};

% ==== 内环反馈 ====
\coordinate (fb2) at ($(G2.east)!0.5!(G1.west)$);
\draw [fline] (fb2) |- ++(0,-1.0) -| (s2.south);
\node at (s2.south) {\tiny $-$};

% ==== 外环反馈 ====
\draw [fline] (out) |- ++(0,-1.5) -| (s1.south);
\node at (s1.south) {\tiny $-$};

% ==== 虚线框 ====
\begin{scope}[on background layer]
    \node [dashbox, fit=(PID1),
           label={[blue!50]above:{\tiny\bfseries 外环}}] {};
    \node [dashbox, fit=(PID2)(G2),
           label={[blue!50]above:{\tiny\bfseries 内环}}] {};
\end{scope}

\end{tikzpicture}
\end{document}
